% ============================================================
%  Глава 3 — AccelPositionEstimator: оценка позиции
%  Файл: pi_nodes/filters/accel_position.py
% ============================================================
\chapter{Оценка позиции по акселерометру с слиянием датчиков}
\label{ch:accel_position}

\begin{filebox}[title={Файл исходного кода}]
\filepath{pi\_nodes/filters/accel\_position.py} \quad (420 строк, Python)
\end{filebox}

\section{Архитектура системы слияния датчиков}

Модуль объединяет три источника информации:

\begin{enumerate}
  \item \textbf{Акселерометр} --- мгновенное линейное ускорение (50 Гц)
  \item \textbf{Колёсная одометрия} --- команды моторов (20 Гц)
  \item \textbf{Orientation from IMU} --- углы Эйлера для поворота фреймов
\end{enumerate}

\section{Константы физики}

\begin{filebox}[title={\filepath{accel\_position.py}, строки 44--69}]
\begin{lstlisting}[style=pythonstyle, numbers=left, firstnumber=44]
EARTH_OMEGA = 7.2921e-5   # rad/s -- Earth angular velocity
DEFAULT_LATITUDE = 55.75  # Moscow

ZUPT_GYRO_ENTER  = 0.015  # rad/s
ZUPT_ACCEL_ENTER = 0.10   # m/s^2
ZUPT_GYRO_EXIT   = 0.04   # rad/s
ZUPT_ACCEL_EXIT  = 0.25   # m/s^2
ZUPT_ENTER_COUNT = 5      # 100ms @ 50Hz
ZUPT_EXIT_COUNT  = 3      # 60ms  @ 50Hz

LATERAL_DECAY = 0.5       # kill 50% lateral velocity per step
\end{lstlisting}
\end{filebox}

\textbf{Угловая скорость Земли} (для компенсации Кориолиса):

\begin{equation}
  \Omega = 7.2921 \times 10^{-5}\,\text{рад/с}
  \label{eq:earth_omega}
\end{equation}

\textbf{Проекция на широту} (Москва, $\varphi = 55.75^\circ$):

\begin{equation}
  \omega_z = \Omega \sin\varphi
  \label{eq:omega_z}
\end{equation}

\section{Шаг 1: Вычитание гравитации}

\begin{filebox}[title={\filepath{accel\_position.py}, строки 211--225 --- update\_imu()}]
\begin{lstlisting}[style=pythonstyle, numbers=left, firstnumber=211]
# Remove gravity -- DIRECT SUBTRACTION
la_bx = ax - self._g_body[0]
la_by = ay - self._g_body[1]
la_bz = az - self._g_body[2]

# Delta orientation correction (slopes only)
d_roll  = roll
d_pitch = pitch
if abs(d_roll) > 0.01 or abs(d_pitch) > 0.01:
    g = self._g_mag
    delta_gx = -g * sin(d_pitch)
    delta_gy  =  g * sin(d_roll)
    la_bx -= delta_gx
    la_by -= delta_gy
\end{lstlisting}
\end{filebox}

\textbf{Прямое вычитание} вектора гравитации, откалиброванного при старте:

\begin{equation}
  \vect{a}^{\text{lin}}_{\text{body}} = \vect{a}_{\text{meas}} - \vect{g}_{\text{body}}
  \label{eq:gravity_subtract}
\end{equation}

\begin{notebox}
\textbf{Почему прямое вычитание, а не ротационный метод?}\\
IMU физически наклонён на шасси. В покое он показывает, например,
$a_x = 0.707\,\text{м/с}^2$ --- чистая гравитация через наклон.
При вычитании: $0.707 - 0.707 = 0$ --- нет ложного движения.
Ротационный метод предполагает $\text{крен}=\text{тангаж}=0$, давая $g_{\text{body}} = (0, 0, 9.81)$,
и оставляет $0.707\,\text{м/с}^2$ как «движение».
\end{notebox}

Поправка на наклон (при нахождении на уклоне $> 0.6^\circ$):

\begin{equation}
  \Delta g_x = -g \sin\theta_{\text{pitch}}, \quad
  \Delta g_y = g \sin\phi_{\text{roll}}
  \label{eq:slope_correction}
\end{equation}

\section{Шаг 2: Поворот из тела в мировую СК}

\begin{filebox}[title={\filepath{accel\_position.py}, строки 227--230}]
\begin{lstlisting}[style=pythonstyle, numbers=left, firstnumber=227]
# Body -> World frame (rotate by yaw)
cy, sy = cos(yaw), sin(yaw)
la_wx = la_bx * cy - la_by * sy
la_wy = la_bx * sy + la_by * cy
\end{lstlisting}
\end{filebox}

Матрица поворота вокруг вертикальной оси ($z$-axis rotation):

\begin{equation}
  \boxed{
  \begin{pmatrix} a^{\text{world}}_x \\ a^{\text{world}}_y \end{pmatrix}
  = \mat{R}(\psi)\,
  \begin{pmatrix} a^{\text{body}}_x \\ a^{\text{body}}_y \end{pmatrix},
  \quad
  \mat{R}(\psi) = \begin{pmatrix} \cos\psi & -\sin\psi \\ \sin\psi & \cos\psi \end{pmatrix}
  }
  \label{eq:body_to_world}
\end{equation}

\section{Шаг 3: Компенсация эффекта Кориолиса}

\begin{filebox}[title={\filepath{accel\_position.py}, строки 232--234}]
\begin{lstlisting}[style=pythonstyle, numbers=left, firstnumber=232]
# Earth rotation compensation (Coriolis)
la_wx -= 2.0 * self._omega_z * self.vy
la_wy += 2.0 * self._omega_z * self.vx
\end{lstlisting}
\end{filebox}

Сила Кориолиса на единицу массы:

\begin{equation}
  \vect{a}_{\text{Cor}} = -2\,\vect{\Omega} \times \vect{v}
  \label{eq:coriolis}
\end{equation}

В 2D (только вертикальная проекция $\omega_z = \Omega\sin\varphi$):

\begin{equation}
  \begin{cases}
    a_{Cx} = -2\,\omega_z\,v_y \\
    a_{Cy} = +2\,\omega_z\,v_x
  \end{cases}
  \label{eq:coriolis_2d}
\end{equation}

Вычитание ускорения Кориолиса из измеренного:

\begin{equation}
  a^{\text{lin}}_x \leftarrow a^{\text{lin}}_x - 2\omega_z v_y, \quad
  a^{\text{lin}}_y \leftarrow a^{\text{lin}}_y + 2\omega_z v_x
  \label{eq:coriolis_remove}
\end{equation}

\section{Шаг 4: Адаптивный медианный фильтр}

\begin{filebox}[title={\filepath{accel\_position.py}, строки 236--250}]
\begin{lstlisting}[style=pythonstyle, numbers=left, firstnumber=236]
# Median filter (reject vibration spikes)
self._buf_ax.append(la_wx)
self._buf_ay.append(la_wy)

if len(self._buf_ax) >= 3:
    la_wx = _median_of_3(
        self._buf_ax[-1], self._buf_ax[-2], self._buf_ax[-3])
    la_wy = _median_of_3(
        self._buf_ay[-1], self._buf_ay[-2], self._buf_ay[-3])

# Noise threshold -- zero out tiny accelerations
if abs(la_wx) < ACCEL_NOISE_THRESHOLD:   # 0.15 m/s^2
    la_wx = 0.0
\end{lstlisting}
\end{filebox}

Медиана трёх последних значений (без сортировки --- одно ветвление):

\begin{equation}
  \hat{a} = \mathrm{median}(a_{k},\, a_{k-1},\, a_{k-2})
  \label{eq:median3}
\end{equation}

Пороговое обнуление шума: $|a| < 0.15\,\text{м/с}^2 \Rightarrow a = 0$.

\section{Шаг 5: ZUPT с гистерезисом}

\begin{filebox}[title={\filepath{accel\_position.py}, строки 252--305}]
\begin{lstlisting}[style=pythonstyle, numbers=left, firstnumber=252]
gyro_mag  = sqrt(gx*gx + gy*gy + gz*gz)
accel_mag = sqrt(la_wx*la_wx + la_wy*la_wy)

if self._stationary:
    # Need STRONG evidence of motion to exit
    if (self._cmd_moving and
        (gyro_mag > ZUPT_GYRO_EXIT or       # 0.04 rad/s
         accel_mag > ZUPT_ACCEL_EXIT)):      # 0.25 m/s^2
        self._exit_count += 1
    if self._exit_count >= ZUPT_EXIT_COUNT:  # 3 samples
        self._stationary = False
else:
    # Easy to re-enter stationary
    if (not self._cmd_moving or
        (gyro_mag < ZUPT_GYRO_ENTER and     # 0.015 rad/s
         accel_mag < ZUPT_ACCEL_ENTER)):    # 0.10 m/s^2
        self._enter_count += 1
    if self._enter_count >= ZUPT_ENTER_COUNT: # 5 samples
        self._stationary = True
\end{lstlisting}
\end{filebox}

\textbf{Двухпороговый автомат состояний (гистерезис):}

\begin{center}
\begin{tikzpicture}[
  state/.style={rectangle, rounded corners, draw=blue!60, fill=blue!10,
                minimum width=3cm, minimum height=1.2cm, font=\bfseries},
  arrow/.style={-{Stealth[length=8pt]}, thick}
]
  \node[state] (S) {STATIONARY};
  \node[state, right=5cm of S] (M) {MOVING};

  \draw[arrow, bend left=20] (S) to node[above, font=\small, align=center]{
    $|\omega|>0.04$ \textbf{или} $|a|>0.25$ м/с$^2$\\
    \textit{за 3 последовательных шага}
  } (M);

  \draw[arrow, bend left=20] (M) to node[below, font=\small, align=center]{
    $|\omega|<0.015$ \textbf{и} $|a|<0.10$ м/с$^2$\\
    \textit{за 5 последовательных шагов}
  } (S);
\end{tikzpicture}
\end{center}

Высокие пороги \textit{выхода} ($> 0.04, > 0.25$) и низкие пороги \textit{входа}
($< 0.015, < 0.10$) предотвращают ложные переходы на границе движения/покоя.

\section{Шаг 6: Трапецеидальное интегрирование (резервный режим)}

\begin{filebox}[title={\filepath{accel\_position.py}, строки 315--334 --- fallback}]
\begin{lstlisting}[style=pythonstyle, numbers=left, firstnumber=315]
# Trapezoidal integration (more accurate than Euler)
self.vx += 0.5 * (la_wx + self._prev_la_x) * dt
self.vy += 0.5 * (la_wy + self._prev_la_y) * dt

# Non-holonomic constraint
fwd = self.vx * cy + self.vy * sy
lat = -self.vx * sy + self.vy * cy
lat *= LATERAL_DECAY
self.vx = fwd * cy - lat * sy
self.vy = fwd * sy + lat * cy

# Position integration
self._accel_x += self.vx * dt
self._accel_y += self.vy * dt
\end{lstlisting}
\end{filebox}

\textbf{Метод трапеций} (порядок точности $O(\Delta t^2)$, лучше Эйлера $O(\Delta t)$):

\begin{equation}
  v_{k} = v_{k-1} + \frac{\Delta t}{2}\bigl(a_{k} + a_{k-1}\bigr)
  \label{eq:trapezoidal}
\end{equation}

\section{Шаг 7: Комплементарный фильтр (резервный режим)}

\begin{filebox}[title={\filepath{accel\_position.py}, строки 376--386 --- blend()}]
\begin{lstlisting}[style=pythonstyle, numbers=left, firstnumber=378]
# Complementary filter: alpha*wheel + (1-alpha)*accel
self._alpha = 0.95 * self._alpha + 0.05 * self._alpha_moving
a = self._alpha
self.x = a * self._wheel_x + (1.0 - a) * self._accel_x
self.y = a * self._wheel_y + (1.0 - a) * self._accel_y

# Slow drift correction of accel toward wheel
self._accel_x = 0.995 * self._accel_x + 0.005 * self._wheel_x
self._accel_y = 0.995 * self._accel_y + 0.005 * self._wheel_y
\end{lstlisting}
\end{filebox}

Комплементарный фильтр с адаптивным $\alpha$:

\begin{equation}
  \vect{p} = \alpha\,\vect{p}_{\text{wheel}} + (1-\alpha)\,\vect{p}_{\text{accel}}
  \label{eq:complementary}
\end{equation}

где $\alpha \in [0.95, 1.0]$ (при движении растёт вес одометрии).
Медленная коррекция дрейфа акселерометра:
$\vect{p}_{\text{accel}} \leftarrow 0.995\,\vect{p}_{\text{accel}} + 0.005\,\vect{p}_{\text{wheel}}$.

\section{Теоретическое обоснование: слияние датчиков}

\subsection{Комплементарный фильтр как частный случай фильтра Калмана}

Комплементарный фильтр~\eqref{eq:complementary} можно рассматривать
как \textbf{стационарный фильтр Калмана} с фиксированным усилением:

\begin{equation}
  \hat{x} = \alpha\,x_{\text{odom}} + (1-\alpha)\,x_{\text{accel}}
  = x_{\text{accel}} + \alpha\,(x_{\text{odom}} - x_{\text{accel}})
  \label{eq:comp_as_kalman}
\end{equation}

Сравнивая с~\eqref{eq:kf_update_x}: $\hat{x} = \hat{x}^- + K\cdot y$,
видим что $\alpha$ играет роль \textbf{усиления Калмана} $K$.

\subsection{Частотная интерпретация}

В частотной области комплементарный фильтр выделяет:
\begin{itemize}
  \item \textbf{Низкие частоты} из одометрии (стабильная, но дрейфующая)
  \item \textbf{Высокие частоты} из акселерометра (шумный, но без дрейфа)
\end{itemize}

\begin{equation}
  X(p) = \underbrace{\frac{\alpha}{1 + \alpha/p}}_{\text{ФНЧ}} X_{\text{odom}}(p)
       + \underbrace{\frac{p/\alpha}{1 + p/\alpha}}_{\text{ФВЧ}} X_{\text{accel}}(p)
  \label{eq:comp_frequency}
\end{equation}

Частота среза: $f_c = \alpha/(2\pi\Delta t) \approx 0.95/(2\pi \cdot 0.02) \approx 7.6\,\text{Гц}$.

\subsection{Трапецеидальное vs.\ Эйлерово интегрирование}

\begin{center}
\begin{tabular}{lll}
\toprule
\textbf{Метод} & \textbf{Формула} & \textbf{Погрешность}\\
\midrule
Эйлер вперёд & $v_k = v_{k-1} + a_k \Delta t$ & $O(\Delta t)$\\
Трапеция & $v_k = v_{k-1} + \frac{a_k + a_{k-1}}{2}\Delta t$ & $O(\Delta t^2)$\\
\bottomrule
\end{tabular}
\end{center}

При частоте $50\,\text{Гц}$ ($\Delta t = 0.02\,\text{с}$) разница в погрешности
составляет $\Delta t/2 = 0.01$, т.е.\ трапеция точнее Эйлера в ${\approx}100$ раз
для гладких сигналов.

\subsection{Дискретная модель пространства состояний}

Полная система AccelPosition в форме пространства состояний
(гл.~\ref{ch:theory}, формула~\eqref{eq:discrete_ss}):

\begin{equation}
  \underbrace{\begin{pmatrix} x \\ y \\ v_x \\ v_y \end{pmatrix}_{k}}_{\vect{x}_k}
  = \underbrace{\begin{pmatrix}
    1 & 0 & \Delta t & 0 \\
    0 & 1 & 0 & \Delta t \\
    0 & 0 & 1 & 0 \\
    0 & 0 & 0 & 1
  \end{pmatrix}}_{\mat{F}}
  \vect{x}_{k-1}
  + \underbrace{\begin{pmatrix}
    \Delta t^2/2 & 0 \\
    0 & \Delta t^2/2 \\
    \Delta t & 0 \\
    0 & \Delta t
  \end{pmatrix}}_{\mat{B}}
  \begin{pmatrix} a_x \\ a_y \end{pmatrix}
  \label{eq:accel_ss}
\end{equation}

\section{Сводная схема конвейера обработки}

\begin{center}
\begin{tikzpicture}[
  node distance=0.6cm,
  block/.style={rectangle, draw=blue!50, fill=blue!8, rounded corners,
                minimum width=4.5cm, minimum height=0.8cm,
                font=\small\bfseries, align=center},
  arrow/.style={-{Stealth}, thick, blue!60}
]
  \node[block] (raw) {Акселерометр \\ $\vect{a}_{\text{meas}}$};
  \node[block, below=of raw] (grav) {Вычитание гравитации \\ $\vect{a}^{\text{lin}} = \vect{a} - \vect{g}_{\text{body}}$};
  \node[block, below=of grav] (rot)  {Поворот СК \\ $\vect{a}^{\text{world}} = \mat{R}(\psi)\vect{a}^{\text{lin}}$};
  \node[block, below=of rot]  (cor)  {Компенсация Кориолиса};
  \node[block, below=of cor]  (med)  {Медианный фильтр + порог шума};
  \node[block, below=of med]  (zupt) {ZUPT с гистерезисом};
  \node[block, below=of zupt] (ekf)  {Velocity EKF \\ (слияние с одометрией)};
  \node[block, below=of ekf]  (nh)   {Неголономное ограничение};
  \node[block, below=of nh]   (pos)  {Позиция $(x, y)$};

  \draw[arrow] (raw)  -- (grav);
  \draw[arrow] (grav) -- (rot);
  \draw[arrow] (rot)  -- (cor);
  \draw[arrow] (cor)  -- (med);
  \draw[arrow] (med)  -- (zupt);
  \draw[arrow] (zupt) -- (ekf);
  \draw[arrow] (ekf)  -- (nh);
  \draw[arrow] (nh)   -- (pos);
\end{tikzpicture}
\end{center}
